\chapter{遥操作的二端口网络与透明度}
\label{ch:teleop-twoport}

% ════════════════════════════════════════════════════════════════
%  机械臂部·遥操作子部第 1 章：二端口网络模型与透明度。
%  定位：把「机器人 + 控制器 + 通信」抽象为电气网络理论的二端口网络，
%  以混合参数 H 统一描述遥操作的因果结构，导出操作者感受阻抗 Z_to 与
%  理想透明度 H_ideal=[[0,1],[-1,0]]，并揭示「透明度 vs 稳定性」的根本矛盾。
%  与 P4 操作空间章 (sec:os-twoport) / 力控章 (sec:fw-twoport) 的二端口骨架互链、
%  不重复；绝对稳定性的完整 Llewellyn 推导、波变量、TDPA 留兄弟章。
%  记号归一：本书统一 P=f^T v 端口功率（承 eq:os-port-power）、tau=J^T f 静力对偶、
%  M(q)q̈+C q̇+g=tau 动力学。源 [F_h;-V_s]=H[V_m;F_e] 已转换到本书黑体记号。
% ════════════════════════════════════════════════════════════════

\begin{practice}[前置自测]
开始之前，先试答下列问题。若已能流畅作答，可只速读本章的因果表与速查卡；若卡住，正是本章要为你解决的。
\begin{enumerate}
  \item 机械阻抗 $Z(s)=f(s)/v(s)$ 与导纳 $Y(s)=v(s)/f(s)$ 各描述什么因果（谁是输入、谁是输出）？一个质量--弹簧--阻尼系统的阻抗写成 $s$ 的函数是什么？
  \item 电路里同一个二端口网络可用 $\mathbf Z$（阻抗）、$\mathbf Y$（导纳）、$\mathbf H$（混合）、$\mathbf T$（传递）四种矩阵描述。它们描述的是同一个系统还是四个不同系统？彼此能互换吗？
  \item 把一台主臂（master）和一台从臂（slave）用电机、传感器、通信线连起来做双边遥操作。操作者希望“像亲手摸到远端环境一样”——这句直觉要求一个怎样的\emph{数学}等式成立？
  \item 一个本身不产生能量的被动反馈网络，接到一个有源元件上，整体一定还是稳定的吗？（想想运放 + RC 网络能搭出振荡器。）
  \item Goertz 在 1950 年代用钢丝绳搭的纯机械主从臂，为什么能同时做到“完美透明”和“绝对稳定”，而后来所有电控方案都做不到？是哪个物理量被电控方案引入了？
\end{enumerate}
\end{practice}

\section*{本章目标}
学完本章，你应当能够：
\begin{enumerate}
  \item \textbf{说清}遥操作从机械联动到电控双边的演化，以及为何“透明度 vs 稳定性”是由物理（通信时延、有源执行）决定的\emph{根本}矛盾而非算法缺陷，并区分 Direct/Shared/Supervisory 三种控制范式；
  \item \textbf{建立}遥操作系统的二端口网络模型：在被动符号约定下定义端口流变量（速度）与势变量（力），并把它与 \cref{ch:operational-space} 的端口功率 $P=\mathbf f^\top\mathbf v$（\cref{eq:os-port-power}）对齐；
  \item \textbf{掌握}混合参数 $\mathbf H$ 的定义与四个元素的物理因果，理解 Hannaford 为何选 $\mathbf H$ 而非 $\mathbf Z/\mathbf Y$（因果匹配 master 力源、slave 运动源），并完成 $\mathbf Z\to\mathbf H$、$\mathbf T\to\mathbf H$ 的转换与级联 $\mathbf T_{\rm tot}=\mathbf T_m\mathbf T_{\rm ch}\mathbf T_s$；
  \item \textbf{推导}操作者感受阻抗 $Z_{to}=h_{11}-h_{12}h_{21}Z_e/(1+h_{22}Z_e)$，并用它解释自由空间/弹簧/刚墙三种环境下操作者“摸到”的等效阻抗；
  \item \textbf{推导}理想透明度的严格条件 $\mathbf H_{\rm ideal}=\bigl[\begin{smallmatrix}0&1\\-1&0\end{smallmatrix}\bigr]$，并逐元素解释其物理含义；
  \item \textbf{对比}四种子架构（PP/PF/FP/FF）的 $\mathbf H$ 矩阵结构与透明度极限，理解它们在工程上的取舍；
  \item \textbf{说明}透明度与绝对稳定性为何在 $\mathbf H_{\rm ideal}$ 处恰好“零裕度”相撞，从而引出兄弟章的绝对稳定性判据（\cref{ch:teleop-stability}）、延迟下的波变量（\cref{ch:teleop-wave}）与时域无源性（\cref{ch:teleop-tdpa}）。
\end{enumerate}

\subsection*{本章知识导航}
本章是一条“\emph{把人--主臂--从臂--环境这条物理链，抽象成一个两端口的黑盒，再用一个 $2\times2$ 矩阵 $\mathbf H$ 量化它有多透明}”的主线。结构如下：

\begin{center}
\small
\begin{tikzpicture}[
  node distance=4mm and 7mm, >=Stealth,
  blk/.style={draw, rounded corners, align=center, font=\small, inner sep=3pt, fill=main!6, draw=main!60!black},
  grp/.style={draw, rounded corners, align=center, font=\small, inner sep=3pt, fill=second!6, draw=second!60!black},
  ln/.style={->, thick, gray!60!black}]
\node[blk] (hist) {核心矛盾\\透明度 vs 稳定性};
\node[blk, right=of hist] (port) {端口变量\\被动符号约定};
\node[blk, right=of port] (h) {混合参数 $\mathbf H$\\因果匹配};
\node[grp, below=8mm of hist] (zto) {感受阻抗 $Z_{to}$\\三特例};
\node[grp, right=of zto] (ideal) {理想透明\\ $\mathbf H_{\rm ideal}$};
\node[grp, right=of ideal] (arch) {四子架构\\ PP/PF/FP/FF};
\node[blk, below=8mm of zto, fill=third!8, draw=third!60!black] (bridge) {桥接 O/P 章\\二端口骨架};
\node[blk, right=of bridge, fill=third!8, draw=third!60!black] (stab) {引出绝对稳定\\\cref{ch:teleop-stability}};
\draw[ln] (hist)--(port); \draw[ln] (port)--(h);
\draw[ln] (h.south) -- ++(0,-4.5mm) -| (zto.north);
\draw[ln] (zto)--(ideal); \draw[ln] (ideal)--(arch);
\draw[ln] (zto.south)--(bridge.north);
\draw[ln] (bridge)--(stab);
\end{tikzpicture}
\end{center}

\noindent\textbf{推荐路径}：动机 $\to$ 端口 $\to$ $\mathbf H$ 三节（\cref{sec:tt-motivation,sec:tt-port,sec:tt-hparam}）是任何读者的主干；感受阻抗 $Z_{to}$ 与理想透明 $\mathbf H_{\rm ideal}$ 两节（\cref{sec:tt-zto,sec:tt-ideal}）是透明度的核心定量工具，后续遥操作各章反复引用；四子架构一节（\cref{sec:tt-arch}）给工程地图；桥接一节（\cref{sec:tt-bridge}）把本章焊接到操作空间/力控的二端口骨架，并把“为什么理想透明不可实现”留给绝对稳定性章。

\subsection*{前置知识桥接}
本章把工具借自两处。其一，\cref{ch:operational-space} 已用\emph{能量端口}视角给出阻抗（运动 $\to$ 力的算子）与导纳（其逆）的二分，并约定端口功率 $P=\mathbf f^\top\mathbf v$（\cref{eq:os-port-power}）、把“机器人 + 控制器”抽象为二端口网络（\cref{sec:os-twoport}，\cref{eq:os-twoport}），还指出“理想透明度要求 $\mathbf H=\bigl[\begin{smallmatrix}\mathbf 0&\mathbf I\\\mathbf I&\mathbf 0\end{smallmatrix}\bigr]$”——本章正是把那条\emph{指针}展开成完整理论。其二，\cref{ch:force-wbc} 的 \cref{sec:fw-twoport} 把四种力控任务统一为二端口网络的不同工作点；本章给出该统一框架在\emph{人在环}遥操作下的母本。此外需要传递函数、Nyquist/BIBO 稳定性的基础（控制理论），以及无源性的端口能量不等式（见 \cref{ch:operational-space} 无源性一节）。

\subsection*{如果跳过本章会怎样}
\begin{itemize}
  \item 在 \cref{ch:teleop-stability}（绝对稳定性）里，你会看到 Llewellyn 四条件直接写在 $h_{11},h_{22},h_{12}h_{21}$ 上，却不知这些符号从何而来、为何 master 用阻抗因果而 slave 用导纳因果——本章正是它们的定义处；
  \item 在 \cref{ch:teleop-wave,ch:teleop-tdpa}（波变量与 TDPA）里，你会发现一切努力都在“用主动控制系统模拟一个被动元件”，而这个目标的\emph{度量}（透明度）与\emph{约束}（被动/稳定）就是本章的 $Z_{to}$ 与 $\mathbf H_{\rm ideal}$。
\end{itemize}

\begin{table}[H]\centering\small
\caption{本章阅读方式建议}
\begin{tabular}{lll}
\toprule
阅读方式 & 时间 & 适合谁 \\
\midrule
精读（含全部推导与练习） & 4--6 小时 & 首次系统学遥操作、要做力反馈系统设计 \\
速读（跳过转换推导细节） & 1.5 小时 & 有二端口基础、想对齐本书记号与因果约定 \\
速查（只看因果表 + 速查卡） & 15 分钟 & 写代码/选架构时回来查 $\mathbf H$ 与 $Z_{to}$ \\
\bottomrule
\end{tabular}
\end{table}

\begin{note}
\textbf{本章记号与约定.}\;
力为势变量（粗体小写或标量 $f$）、速度为流变量 $v$；机械阻抗 $Z(s)=f/v$、导纳 $Y(s)=v/f$。端口功率沿用全书 $P=\mathbf f^\top\mathbf v$（\cref{eq:os-port-power}），\emph{流入网络}的速度取正。主臂侧物理量记 $f_h$（人手施于 master 的力）、$v_m$（master 速度）；从臂侧 $f_e$（环境施于 slave 的力）、$v_s$（slave 速度）；环境阻抗 $Z_e$、操作者阻抗 $Z_h$。
\textbf{源文献}（Hannaford 1989、Lawrence 1993）把混合矩阵写成 $[F_h;-V_s]^\top=\mathbf H[V_m;F_e]^\top$、并以 $s$ 域大写表频域量；本书统一为下式 \cref{eq:tt-hdef} 的形式（端口 2 速度取 $-v_s$ 以使两端都用“流入网络”的流变量记账，与 \cref{eq:os-twoport} 的 $\mathbf H$ 一致）。运动/力缩放（$x_s=x_m/\lambda$、$f_h=\lambda f_e$）等若引用源约定，本书转换后在 note 中注明。
\end{note}

% ════════════════════════════════════════════════════════════════
\section{动机：从完美的钢丝绳到会振荡的电控}
\label{sec:tt-motivation}

\paragraph{动机：人要在够不到的地方动手。}
遥操作（teleoperation）的需求朴素而古老：在放射、深海、太空、微创手术这些人\emph{进不去}或\emph{不该进去}的环境里完成物理操作。把这件事做到极致，意味着操作者既要能\emph{精确地动}（运动从主臂传到从臂），又要能\emph{真切地感}（接触力从从臂传回主臂手上）。这后半句——让操作者感受到远端的真实力——正是本章的主角“透明度”。

\paragraph{反面：电控一上来就既不稳又不真。}
第一代遥操作是\emph{纯机械}的。Goertz 在阿贡国家实验室（机械主从臂起于 1940 年代末——金属带式 1948、专利 1949、钢索滑轮式约 1951；其后才发展出 1954 年的电控/伺服力反射版）用钢丝绳与滑轮把隔离墙两侧的主、从臂直接连起来处理放射性材料\cite{goertz1954mechanical}。这套\emph{机械}系统有一个后来再没被电控方案完整复现的优良性质：

\begin{insight}[机械主从臂为何“完美”]\label{ins:tt-goertz}
钢丝绳连接是一个\emph{纯被动}机械系统——没有电机、没有控制器、没有通信延迟。它不可能凭空产生能量，故\textbf{绝对稳定}；而刚性传动使 $f_h\approx f_e$、$v_m\approx v_s$（理想刚体），故\textbf{完美透明}。后续一切理论（二端口网络、波变量、TDPA）本质上都是在\emph{用主动控制系统去模拟这个被动元件的行为}。透明与稳定之所以在机械方案里“免费”同时成立，根因是：信号沿钢丝以声速（约 $5000\,\mathrm{m/s}$）传播，$3\,\mathrm{m}$ 距离延迟仅约 $0.6\,\mathrm{ms}$，对人类感知不可察觉。
\end{insight}

机械联动的死穴是\emph{距离}（约 $3\,\mathrm{m}$）与无缩放。第二代电控方案（1960--1980 年代）用电机替代钢丝绳、用传感器与通信线替代机械连接，立刻解锁了任意距离与运动/力缩放——却同时引入三个新问题：执行器惯量与齿轮摩擦让操作者\emph{感受不到}真实环境（透明度下降）；控制器延迟与离散采样让系统在硬接触时\emph{振荡}（稳定性下降）；而最致命的是\textbf{通信延迟}。Ferrell 在 MIT 的实验（1965）首次定量揭示了延迟的摧毁性\cite{ferrell1965remote}：当往返延迟进入人类反应时尺度（约数百毫秒）以上，人类便从“连续控制”退化为“动一下--等一等--看一眼”（move-and-wait）的策略，且任务完成时间随延迟近似\emph{线性}增长。\pz{“切换为 move-and-wait”的精确延迟阈值依任务/反馈方式而异：Ferrell 1965 原实验所加延迟可达秒级，文献给出的明显退化阈值从约数百毫秒到 1--2 秒不等，本处“约数百毫秒”取人类反应时量级作定性界，非单一硬阈值}

\begin{figure}[ht]
\centering
\begin{tikzpicture}[scale=1.0,>=Stealth, font=\small]
  % axes
  \draw[->] (0,0)--(6.6,0) node[right]{延迟 $T$};
  \draw[->] (0,0)--(0,4.2) node[above]{任务时间倍数};
  \foreach \y/\l in {1/1$\times$,2/2$\times$,3/3$\times$,4/4$\times$}{
    \draw (-0.06,\y)--(0.06,\y); \node[left,scale=0.8] at (-0.08,\y){\l};}
  \foreach \x/\l in {0.4/0,1.6/50,3.0/200,4.4/500,5.8/1000}{
    \draw (\x,-0.06)--(\x,0.06); \node[below,scale=0.75] at (\x,-0.08){\l\,ms};}
  % data curve (schematic, after Ferrell 1965)
  \draw[very thick, main] (0.4,1.0) .. controls (1.0,1.05) and (1.4,1.15) .. (1.6,1.2)
        .. controls (2.3,1.5) and (2.7,1.85) .. (3.0,2.0)
        .. controls (3.7,2.6) and (4.1,3.2) .. (4.4,3.5)
        .. controls (5.0,3.9) and (5.4,4.0) .. (5.8,4.05);
  \fill[main] (0.4,1.0) circle(1.4pt) (1.6,1.2) circle(1.4pt) (3.0,2.0) circle(1.4pt)
              (4.4,3.5) circle(1.4pt);
  % move-and-wait threshold
  \draw[dashed, red!60!black] (3.0,0)--(3.0,2.0);
  \node[red!60!black, align=center, scale=0.8] at (4.6,1.15)
       {反应时尺度以上\\渐入 move-and-wait};
\end{tikzpicture}
\caption{Ferrell（1965）时延实验的示意趋势：任务完成时间随往返延迟近似线性增长，越过人类反应时尺度后操作策略逐渐转为 move-and-wait（图中虚线为示意位置，非精确阈值）。曲线为示意重绘，强调趋势而非精确数值。\pz{存疑：各延迟点的具体时间倍数与策略切换的精确阈值源自二手转述，需以 Ferrell 1965 原文核实；原实验所加延迟可达秒级}}
\label{fig:tt-ferrell}
\end{figure}

\paragraph{核心思想：透明度与稳定性是一对结构性对手。}
本章要反复打磨的一句话是：

\begin{insight}[遥操作的核心矛盾]\label{ins:tt-core}
\emph{操作者越想完美感受远端环境（透明度越高），系统越容易因延迟与有源执行而失稳。}这不是“换个更好的算法”能消除的——它是物理定律（能量守恒、因果性、信息传播有限速）施加的根本权衡。我们能做的不是消灭矛盾，而是\emph{量化}它、并在两端之间找最佳折中。本章建立量化它的数学工具：二端口网络与透明度；下一章（\cref{ch:teleop-stability}）给出稳定性这一侧的精确边界。
\end{insight}

\paragraph{三种控制范式：人和机器谁做主。}
在展开 Direct 范式的精细分析前，先把全景摆正。按人与机器人之间\emph{控制权}如何分配，遥操作分三类（\cref{tab:tt-paradigm}）。本章及整个遥操作子部聚焦 \textbf{Direct（直接遥操作）}——操作者连续控制每个自由度、机器人提供力反馈——因为它对透明度要求最高、理论分析最深。

\begin{table}[H]\centering\small
\caption{遥操作的三种控制范式}
\label{tab:tt-paradigm}
\begin{tabular}{llll}
\toprule
范式 & 人的角色 & 机器自主度 & 典型应用 \\
\midrule
Direct（直接） & 连续控制所有 DOF & 零（纯执行） & 核处理、爆炸物处理 \\
Shared（共享） & 提供高级意图 & 部分（自动避障等） & 手术机器人（da Vinci） \\
Supervisory（监督） & 下达目标 / 监控 & 高（自主规划执行） & 深空探测（火星车） \\
\bottomrule
\end{tabular}
\end{table}

\begin{note}
\textbf{跨领域类比.}\;三种范式恰似自动驾驶的分级：Direct $\approx$ L0（人全控）、Shared $\approx$ L2--L3（人机共驾）、Supervisory $\approx$ L4--L5（机器自主）。一如自动驾驶中 L3 的“人机切换”最危险，遥操作里 Shared 范式的“人机意图冲突”也最棘手——这部分留待 \cref{ch:teleop-mapping} 的运动映射与共享自主。
\end{note}

% ════════════════════════════════════════════════════════════════
\section{端口变量与被动符号约定}
\label{sec:tt-port}

\paragraph{直觉：把整条链塞进一个黑盒，只看它两端。}
人--主臂--通信--从臂这条链内部千头万绪（电机、编码器、控制器、延迟）。二端口网络（two-port network）的智慧在于：\emph{不看内部}，只把它当一个黑盒，黑盒有两个端口，端口 1 接操作者、端口 2 接环境（\cref{fig:tt-twoport}）。每个端口由一对变量刻画——一个\emph{势变量}（力 $f$）、一个\emph{流变量}（速度 $v$），二者乘积是\emph{功率} $P=fv$。这套语言搬自微波电路理论，由 Hannaford（1989）引入遥操作\cite{hannaford1989design}，与 \cref{ch:operational-space} 把“机器人 + 控制器”抽象为二端口（\cref{sec:os-twoport}）是同一思想在\emph{人在环}场景下的延伸。

\begin{table}[H]\centering\small
\caption{电气--力学对偶（沿用 \cref{eq:os-port-power} 的端口功率约定）}
\label{tab:tt-analogy}
\begin{tabular}{llll}
\toprule
电气量 & 力学量 & 符号 & 单位 \\
\midrule
电压 $V$ & 力 $f$ & $f$ & N \\
电流 $I$ & 速度 $v$ & $v$ & m/s \\
功率 $P=VI$ & 功率 $P=fv$ & $P$ & W \\
阻抗 $Z=V/I$ & 阻抗 $Z=f/v$ & $Z$ & N\,s/m \\
\bottomrule
\end{tabular}
\end{table}

\begin{figure}[ht]
\centering
\begin{tikzpicture}[>=Stealth, font=\small]
  \node[draw, rounded corners, minimum width=3.4cm, minimum height=2.2cm,
        fill=main!6, draw=main!60!black, align=center] (box) {遥操作系统\\（黑盒 $\mathbf H$）};
  % port 1 labels (left)
  \node[left=1.9cm of box.west, align=center] (op) {操作者\\端口 1};
  \draw[->, thick] (op.east) ++(0,0.45) -- node[above,scale=0.85]{$f_h$} ([yshift=0.45cm]box.west);
  \draw[<-, thick] (op.east) ++(0,-0.45) -- node[below,scale=0.85]{$v_m$} ([yshift=-0.45cm]box.west);
  % port 2 labels (right)
  \node[right=1.9cm of box.east, align=center] (env) {环境\\端口 2};
  \draw[<-, thick] (env.west) ++(0,0.45) -- node[above,scale=0.85]{$f_e$} ([yshift=0.45cm]box.east);
  \draw[->, thick] (env.west) ++(0,-0.45) -- node[below,scale=0.85]{$v_s$} ([yshift=-0.45cm]box.east);
  % passive sign note
  \node[below=0.2cm of box, align=center, scale=0.82, gray!50!black]
    {流入网络的流变量取正：端口 1 用 $\nu_1=v_m$，端口 2 用 $\nu_2=-v_s$};
\end{tikzpicture}
\caption{遥操作二端口网络拓扑。$f_h$、$f_e$ 指向网络内部（操作者推 master、环境推 slave）；$v_m,v_s$ 为物理运动正方向。做无源性记账时，两端都用“流入网络”的流变量，故端口 2 取 $\nu_2=-v_s$。}
\label{fig:tt-twoport}
\end{figure}

\paragraph{被动符号约定：为什么端口 2 的速度要带负号。}
能量记账要求把每个端口的功率写成统一形式 $P_k=f_k\,\nu_k$，其中 $\nu_k$ 是\emph{流入}网络的流变量。端口 1 处，操作者推 master、master 以 $v_m$ 运动，流入网络的流就是 $\nu_1=v_m$。端口 2 处，环境施力 $f_e$ 指向网络内、而 slave 以 $v_s$ \emph{离开}网络运动，故流入网络的流是 $\nu_2=-v_s$。于是网络吸收的总功率

\begin{equation}
  P = f_h\,v_m + f_e\,(-v_s) = f_h v_m - f_e v_s,
  \label{eq:tt-port-power}
\end{equation}

与 \cref{eq:os-port-power} 的 $P=\mathbf f^\top\mathbf v$ 在每个端口上逐项对应。这一处看似琐碎的负号，是后文一切无源性/绝对稳定性分析的记账基准；若两端都用物理速度 $v_s$（不取负），无源性不等式会系统性地差一个符号。把它一次性钉死，后面就再不必纠结。

\begin{pitfall}[符号陷阱：物理方向 vs 记账方向]
新手常把 $v_s$ 的“物理正方向”与“功率记账方向”混为一谈。物理上 $v_s$ 当然是 slave 真实运动的方向，画图标注没错；但\emph{做能量/无源性记账}时，端口 2 的流变量必须取 $\nu_2=-v_s$，否则 $P=f_hv_m+f_ev_s$ 会把“从端口 2 流出的功率”误记为“流入”。本章后续的 $\mathbf H$ 矩阵之所以把输出写成 $-v_s$（见 \cref{eq:tt-hdef}），正是为了让矩阵直接吃/吐\emph{记账方向}的量。
\end{pitfall}

% ════════════════════════════════════════════════════════════════
\section{混合参数 H：遥操作的因果母本}
\label{sec:tt-hparam}

\subsection{定义与四元素的物理因果}
\label{sec:tt-hdef}

线性二端口网络的内部关系由四个独立参数决定。\emph{选哪两个量作自变量}，就得到不同的参数化（\cref{ch:operational-space} 的 \cref{eq:os-twoport} 已列出 $\mathbf Z/\mathbf Y/\mathbf H$ 三种）。遥操作选\textbf{混合参数（hybrid parameters）}$\mathbf H$：

\begin{definition}[混合参数矩阵]\label{def:tt-hmatrix}
遥操作二端口网络的混合参数矩阵 $\mathbf H(s)$ 定义为
\begin{equation}
  \begin{bmatrix} f_h(s) \\[1pt] -v_s(s) \end{bmatrix}
  =\underbrace{\begin{bmatrix} h_{11}(s) & h_{12}(s) \\[1pt] h_{21}(s) & h_{22}(s) \end{bmatrix}}_{\mathbf H(s)}
  \begin{bmatrix} v_m(s) \\[1pt] f_e(s) \end{bmatrix}.
  \label{eq:tt-hdef}
\end{equation}
即以\emph{端口 1 的速度} $v_m$ 与\emph{端口 2 的力} $f_e$ 为输入，以端口 1 的力 $f_h$ 与端口 2 的（记账方向）速度 $-v_s$ 为输出。
\end{definition}

每个 $h_{ij}$ 都有干净的物理意义，可由“开路/短路”测量读出（\cref{tab:tt-hphys}）。这里“环境断开”指 $f_e=0$（从臂悬空、自由空间），“master 固定”指 $v_m=0$（夹死主臂）。

\begin{table}[H]\centering\small
\caption{$\mathbf H$ 四元素的测量条件与物理因果}
\label{tab:tt-hphys}
\begin{tabular}{lll}
\toprule
参数 & 测量条件 & 物理含义 \\
\midrule
$h_{11}=\dfrac{f_h}{v_m}\Big|_{f_e=0}$ & 环境断开（自由空间） & master 自由空间\emph{阻抗}（惯量/摩擦） \\[6pt]
$h_{12}=\dfrac{f_h}{f_e}\Big|_{v_m=0}$ & master 固定 & 力反射比（环境力 $\to$ 操作者手感） \\[6pt]
$h_{21}=\dfrac{-v_s}{v_m}\Big|_{f_e=0}$ & 环境断开 & 运动传递比（位置跟踪） \\[6pt]
$h_{22}=\dfrac{-v_s}{f_e}\Big|_{v_m=0}$ & master 固定 & slave \emph{导纳}（环境力 $\to$ slave 让步） \\
\bottomrule
\end{tabular}
\end{table}

\paragraph{为什么选 $\mathbf H$ 而非 $\mathbf Z$ 或 $\mathbf Y$。}
关键不是数学方便，而是\emph{因果匹配}。$\mathbf Z$ 参数 $[f_1;f_2]=\mathbf Z[v_1;v_2]$ 让两端都“输入速度、输出力”（两端皆阻抗因果）；$\mathbf Y$ 参数 $[v_1;v_2]=\mathbf Y[f_1;f_2]$ 让两端都“输入力、输出速度”（两端皆导纳因果）。但真实遥操作里，\emph{master 端是力源}（操作者用力推，系统反馈力让你感受）、\emph{slave 端是运动源}（跟随主臂运动，环境力使其让步）。$\mathbf H$ 恰好“混合”这两种因果：端口 1 阻抗因果（$v_m\to f_h$）、端口 2 导纳因果（$f_e\to v_s$）。这正是 \cref{ch:operational-space} 反复强调的“互连两端不能都是导纳/都是阻抗”在遥操作里的具体落点。

\begin{pitfall}[纠正：“遥操作必须从散射参数学起”]
不少教材从散射参数 $\mathbf S$ 或 $\mathbf Z$ 切入遥操作。但 Hannaford 选 $\mathbf H$ 不是为了数学优雅，而是因为 $\mathbf H$ \emph{直接对应物理因果结构}——master 力源、slave 运动源。散射参数 $\mathbf S$ 的长处在别处：它把无源性写成 $\|\mathbf S\|_\infty\le1$ 的简洁形式，故在\emph{延迟下的无源通信}里称王（\cref{ch:teleop-wave} 的波变量正是 $\mathbf S$ 的舞台）。同一系统、不同参数化，按分析目的选用，彼此无损互换。
\end{pitfall}

\subsection{从 Z 与 T 转换到 H：一条参数化链}
\label{sec:tt-convert}

五种参数化（$\mathbf Z,\mathbf Y,\mathbf H,\mathbf T,\mathbf S$）描述的是\emph{同一个}系统，转换不丢信息。下面给两条最常用的转换，它们既是练手，也揭示各参数化的分工。

\begin{derivation}[$\mathbf Z\to\mathbf H$]
从阻抗参数 $f_1=z_{11}v_1+z_{12}v_2$、$f_2=z_{21}v_1+z_{22}v_2$ 出发，改以 $(v_1,f_2)$ 为自变量。第二式解出 $v_2=(f_2-z_{21}v_1)/z_{22}$，代入第一式：
\begin{equation}
  f_1=\Big(z_{11}-\tfrac{z_{12}z_{21}}{z_{22}}\Big)v_1+\tfrac{z_{12}}{z_{22}}f_2 .
\end{equation}
注意 \cref{eq:tt-hdef} 的输出第二行是 $-v_2$，故对 $v_2$ 表达式取负：$-v_2=\tfrac{z_{21}}{z_{22}}v_1-\tfrac{1}{z_{22}}f_2$。读出系数：
\begin{equation}
  h_{11}=z_{11}-\frac{z_{12}z_{21}}{z_{22}},\quad
  h_{12}=\frac{z_{12}}{z_{22}},\quad
  h_{21}=\frac{z_{21}}{z_{22}},\quad
  h_{22}=-\frac{1}{z_{22}}.
  \label{eq:tt-z2h}
\end{equation}
\end{derivation}

传递矩阵 $\mathbf T$（ABCD 参数）单独值得一提，因为它有一个其它参数化没有的优良性质——\emph{级联即相乘}。

\begin{proposition}[级联性与 $\mathbf T\to\mathbf H$]\label{prop:tt-cascade}
以 $(f_2,v_2)$ 表 $(f_1,v_1)$ 的传递矩阵
\begin{equation}
  \begin{bmatrix}f_1\\v_1\end{bmatrix}
  =\begin{bmatrix}A&B\\C&D\end{bmatrix}\begin{bmatrix}f_2\\v_2\end{bmatrix}
  \label{eq:tt-tdef}
\end{equation}
满足：两个二端口网络串联时总传递矩阵等于因子之积。对遥操作整链，
\begin{equation}
  \mathbf T_{\rm tot}=\mathbf T_{\rm master}\,\mathbf T_{\rm channel}\,\mathbf T_{\rm slave},
  \label{eq:tt-Ttot}
\end{equation}
其中 $\mathbf T_{\rm channel}$ 编码通信通道（含延迟 $e^{-sT}$）。由 \cref{eq:tt-tdef} 改以 $(v_1,f_2)$ 为自变量（第二式解 $v_2=(v_1-Cf_2)/D$ 代入第一式），得
\begin{equation}
  h_{11}=\frac{B}{D},\quad
  h_{12}=\frac{AD-BC}{D}=\frac{\det\mathbf T}{D},\quad
  h_{21}=-\frac{1}{D},\quad
  h_{22}=\frac{C}{D}.
  \label{eq:tt-t2h}
\end{equation}
\end{proposition}
\begin{proof}
级联性：设网络 $\mathcal N_1$ 的输出端口 $(f_2^{(1)},v_2^{(1)})$ 直接接 $\mathcal N_2$ 的输入端口 $(f_1^{(2)},v_1^{(2)})$，物理连接给 $f_2^{(1)}=f_1^{(2)}$、$v_2^{(1)}=v_1^{(2)}$（同一界面力与速度连续）。于是 $[f_1^{(1)};v_1^{(1)}]=\mathbf T_1[f_2^{(1)};v_2^{(1)}]=\mathbf T_1\mathbf T_2[f_2^{(2)};v_2^{(2)}]$，即总传递矩阵为 $\mathbf T_1\mathbf T_2$。\cref{eq:tt-t2h} 的代数整理同 $\mathbf Z\to\mathbf H$，从略。
\end{proof}

\begin{insight}[各参数化的分工]\label{ins:tt-param-roles}
$\mathbf Z/\mathbf Y$ 宜做\emph{阻抗/导纳分析}（环境耦合、力控架构）；$\mathbf T$ 宜做\emph{级联}（master--channel--slave 串联，\cref{eq:tt-Ttot}）；$\mathbf H$ 宜做\emph{遥操作标准描述}（因果匹配，本章主角）；$\mathbf S$ 宜做\emph{无源性/延迟分析}（\cref{ch:teleop-wave}）。理解它们互为线性分式变换、可无损互换，就握住了遥操作理论的“统一语言”。$\mathbf H\leftrightarrow\mathbf S$ 的线性分式变换（$\mathbf S=(\mathbf C+\mathbf D\mathbf H)(\mathbf A+\mathbf B\mathbf H)^{-1}$ 形式，系数由波阻抗 $b$ 定）将在 \cref{ch:teleop-wave} 展开。
\end{insight}

% ════════════════════════════════════════════════════════════════
\section{操作者感受阻抗 \texorpdfstring{$Z_{to}$}{Zto}}
\label{sec:tt-zto}

透明度需要一个\emph{可计算的度量}。这个度量就是操作者感受阻抗 $Z_{to}$——操作者通过整个遥操作系统“摸到”的等效阻抗。它把抽象的“手感”翻译成一个传递函数。

\begin{theorem}[操作者感受阻抗]\label{thm:tt-zto}
设环境呈线性阻抗 $f_e=Z_e\,v_s$。接到 \cref{eq:tt-hdef} 的端口 2 后，操作者在 master 端看到的等效阻抗 $Z_{to}:=f_h/v_m$ 为
\begin{equation}
  Z_{to}(s)=h_{11}-\frac{h_{12}h_{21}\,Z_e}{1+h_{22}Z_e}
  =\frac{h_{11}+\Delta_h\,Z_e}{1+h_{22}Z_e},
  \qquad \Delta_h:=h_{11}h_{22}-h_{12}h_{21},
  \label{eq:tt-zto}
\end{equation}
其中 $\Delta_h=\det\mathbf H$。
\end{theorem}
\begin{proof}
把 $f_e=Z_e v_s$ 代入 \cref{eq:tt-hdef} 第二行：$-v_s=h_{21}v_m+h_{22}Z_e v_s$，解得
\begin{equation}
  v_s=\frac{-h_{21}v_m}{1+h_{22}Z_e},\qquad
  f_e=Z_e v_s=\frac{-h_{21}Z_e v_m}{1+h_{22}Z_e}.
\end{equation}
代入第一行 $f_h=h_{11}v_m+h_{12}f_e$：
\begin{equation}
  f_h=h_{11}v_m+h_{12}\cdot\frac{-h_{21}Z_e v_m}{1+h_{22}Z_e}
     =\Big(h_{11}-\frac{h_{12}h_{21}Z_e}{1+h_{22}Z_e}\Big)v_m,
\end{equation}
两边除以 $v_m$ 即 \cref{eq:tt-zto} 的第一式；通分并用 $\Delta_h$ 的定义得第二式。
\end{proof}

\paragraph{逐项的物理解读。}
$h_{11}$ 是\emph{环境无关}的本底——即便自由空间（$Z_e=0$），操作者也会感到 master 自身的惯量与摩擦，$Z_{to}=h_{11}$。第二项 $\tfrac{h_{12}h_{21}Z_e}{1+h_{22}Z_e}$ 是环境阻抗 $Z_e$ 经力反射（$h_{12}$）与运动传递（$h_{21}$）“穿透”系统后的等效值；分母里的 $h_{22}Z_e$ 表征 slave 导纳对穿透的衰减。$\Delta_h=\det\mathbf H$ 则决定高频（$Z_e\to\infty$）时的传递增益。三种典型环境的极限给出最直观的体检（\cref{tab:tt-zto-cases}）。

\begin{table}[H]\centering\small
\caption{三种典型环境下的 $Z_{to}$ 与理想值}
\label{tab:tt-zto-cases}
\begin{tabular}{llll}
\toprule
环境 & $Z_e$ & $Z_{to}$（\cref{eq:tt-zto}） & 理想值 \\
\midrule
自由空间 & $0$ & $h_{11}$ & $0$ \\[2pt]
纯弹簧 & $K/s$ & $h_{11}-\dfrac{h_{12}h_{21}K}{s+h_{22}K}$ & $K/s$ \\[6pt]
刚墙 & $\infty$ & $\dfrac{\Delta_h}{h_{22}}$ & $\infty$ \\
\bottomrule
\end{tabular}
\end{table}

\begin{pitfall}[概念误区：“力和速度都跟得上 = 透明”]
直觉常以为 $f_h=f_e$ 且 $v_m=v_s$ 就是完美透明。但仅信号匹配不够——透明度定义在\emph{阻抗关系}上。设想一个缩放系统 $f_h=\lambda f_e$、$v_s=\lambda v_m$：力与速度都“不等”，可操作者感到的 $Z_{to}=f_h/v_m=\lambda f_e/(v_s/\lambda)=\lambda^2 f_e/v_s$，是否透明取决于缩放是否一致地作用于阻抗。\textbf{透明度的本质是 $Z_{to}=Z_e$ 这一比信号匹配更深的条件}——这也是下一节为何用 $Z_{to}$ 而非力跟踪误差来定义它。
\end{pitfall}

% ════════════════════════════════════════════════════════════════
\section{理想透明度与 \texorpdfstring{$\mathbf H_{\rm ideal}$}{Hideal}}
\label{sec:tt-ideal}

\begin{definition}[理想透明度，Lawrence 1993]\label{def:tt-transparency}
遥操作系统是\emph{理想透明}的，当且仅当操作者感受阻抗恒等于环境阻抗：
\begin{equation}
  Z_{to}(s)\equiv Z_e(s),\qquad \forall Z_e,\ \forall s.
  \label{eq:tt-transparency}
\end{equation}
即对\emph{任意}环境、\emph{所有}频率，操作者感到的都正是真实环境阻抗——系统“消失”了，如同亲手触碰\cite{lawrence1993stability}。
\end{definition}

\paragraph{从 $Z_{to}\equiv Z_e$ 反解 $\mathbf H$。}
理想透明度对 $\mathbf H$ 的约束极强，强到几乎唯一地钉死它。

\begin{theorem}[理想透明的 $\mathbf H$ 矩阵]\label{thm:tt-hideal}
满足 \cref{eq:tt-transparency} 的混合参数矩阵为
\begin{equation}
  \mathbf H_{\rm ideal}=\begin{bmatrix} 0 & 1 \\ -1 & 0 \end{bmatrix}.
  \label{eq:tt-hideal}
\end{equation}
\end{theorem}
\begin{proof}
令 \cref{eq:tt-zto} 的第二式等于 $Z_e$：$\dfrac{h_{11}+\Delta_h Z_e}{1+h_{22}Z_e}=Z_e$，即
\begin{equation}
  h_{11}+\Delta_h Z_e = Z_e + h_{22}Z_e^2.
\end{equation}
这须对\emph{所有} $Z_e$ 成立，故各次幂系数分别相等（多项式恒等定理）：
\begin{equation}
  Z_e^0:\ h_{11}=0;\qquad
  Z_e^1:\ \Delta_h=1;\qquad
  Z_e^2:\ h_{22}=0.
\end{equation}
由 $h_{11}=h_{22}=0$，行列式 $\Delta_h=h_{11}h_{22}-h_{12}h_{21}=-h_{12}h_{21}=1$，即 $h_{12}h_{21}=-1$。最简（且对称）解为 $h_{12}=1,\ h_{21}=-1$，得 \cref{eq:tt-hideal}。代回 \cref{eq:tt-zto} 验证：$Z_{to}=0-\dfrac{1\cdot(-1)\cdot Z_e}{1+0}=Z_e$ $\checkmark$。
\end{proof}

\paragraph{四个元素的物理验证。}
\cref{eq:tt-hideal} 的每个元素都说得通：$h_{11}=0$ 表示自由空间 master 阻抗为零——操作者感受不到 master 自身的惯量与摩擦；$h_{22}=0$ 表示 slave 导纳为零——master 固定时 slave 完全不被环境力推动（刚性运动源）；$h_{12}=1$ 表示环境力 $100\%$ 传回操作者；$h_{21}=-1$ 表示 master 速度 $100\%$ 传到 slave（负号是 \cref{eq:tt-hdef} 输出取 $-v_s$ 的记账约定）。这与 \cref{ch:operational-space} 在 \cref{eq:os-twoport} 之后给出的 $\mathbf H=\bigl[\begin{smallmatrix}\mathbf 0&\mathbf I\\\mathbf I&\mathbf 0\end{smallmatrix}\bigr]$ 标量版完全一致（该处 $\mathbf I$ 块对应这里的 $\pm1$，符号差源于端口 2 速度记账约定）。

\begin{insight}[透明度的物理本质]\label{ins:tt-transparency-essence}
理想透明要求 $h_{11}=h_{22}=0$——即 master 与 slave 的“自身存在”完全消失。但任何物理实体都有惯量、摩擦、延迟，故实际系统 $h_{11},h_{22}\neq0$。\textbf{透明度的本质，是让遥操作系统在操作者看来“不存在”；而任何物理装置都不可能真正不存在}——这是透明度永远只能逼近、不能达到的物理根由，也是 \cref{sec:tt-bridge} 将揭示的“透明 vs 稳定”相撞的伏笔。
\end{insight}

% ════════════════════════════════════════════════════════════════
\section{四种子架构：PP / PF / FP / FF}
\label{sec:tt-arch}

Hannaford（1989）把遥操作控制器分解为四个独立通道（\cref{fig:tt-fourchan}）：位置前馈 $C_1$（master 速度 $\to$ slave 指令）、位置反馈 $C_2$（slave 速度 $\to$ master）、力前馈 $C_3$（master 力 $\to$ slave）、力反馈 $C_4$（slave 力 $\to$ master）。选择性开关这四个通道，得到四种“子架构”，以两端各自传递的是位置（P）还是力（F）命名。

\begin{figure}[ht]
\centering
\begin{tikzpicture}[>=Stealth, font=\small, node distance=6mm]
  \node[draw, rounded corners, minimum width=2.3cm, minimum height=1.5cm,
        fill=second!8, draw=second!60!black, align=center] (m) {master\\控制器};
  \node[draw, rounded corners, minimum width=2.3cm, minimum height=1.5cm,
        fill=third!8, draw=third!60!black, align=center, right=4.3cm of m] (s) {slave\\控制器};
  % C1 position forward (top, m->s)
  \draw[->, thick, main!60!black] (m.north east) ++(0,-0.15)
      to[out=20,in=160] node[above,scale=0.82]{$C_1$ 位置 m$\to$s} ([yshift=0.35cm]s.north west);
  % C3 force forward (m->s) lower
  \draw[->, thick, main!60!black] (m.east) ++(0,0.25)
      to[out=8,in=172] node[above,scale=0.82,pos=0.5]{$C_3$ 力 m$\to$s} ([yshift=0.1cm]s.west);
  % C4 force feedback (s->m)
  \draw[->, thick, red!55!black] (s.west) ++(0,-0.25)
      to[out=188,in=-8] node[below,scale=0.82,pos=0.5]{$C_4$ 力 s$\to$m} ([yshift=-0.1cm]m.east);
  % C2 position feedback (s->m) bottom
  \draw[->, thick, red!55!black] (s.south west) ++(0,0.15)
      to[out=200,in=-20] node[below,scale=0.82]{$C_2$ 位置 s$\to$m} ([yshift=-0.35cm]m.south east);
\end{tikzpicture}
\caption{Hannaford 四通道架构。$C_1,C_3$ 为前馈（master$\to$slave），$C_2,C_4$ 为反馈（slave$\to$master）。开关组合给出 PP/PF/FP/FF 四子架构与“理想四通道”全开。}
\label{fig:tt-fourchan}
\end{figure}

\paragraph{四子架构的 H 矩阵结构与透明度极限。}
逐架构的完整 $\mathbf H$ 推导冗长且依赖具体控制器，本章只给\emph{结构性}对比（\cref{tab:tt-arch}）——透明度由 $h_{11},h_{22}$ 离零多远、$h_{12}h_{21}$ 离 $-1$ 多远决定（对照 \cref{eq:tt-hideal}）。一个反复出现的纪律：

\begin{pitfall}[记号陷阱：$\mathbf H$ 里不许出现 $Z_e$]
有的资料把某架构的“$\mathbf H$ 矩阵”写成含 $Z_e$ 的式子。这是错的：$\mathbf H(s)$ 描述的是遥操作\emph{二端口本体}，环境阻抗 $Z_e$ 是\emph{接到端口 2 的终端}，只能在算 $Z_{to}$（\cref{eq:tt-zto}）时以 $f_e=Z_e v_s$ 代入才出现。若一个“$\mathbf H$”里含 $Z_e$，它描述的已是“接上特定环境后的闭环”，而非二端口参数本身。判别法：$\mathbf H$ 的四元素必须只由 master/slave 本体动力学、控制器、通信环节决定。
\end{pitfall}

\begin{itemize}
  \item \textbf{PP（位置--位置）}：只用 $C_1,C_2$，$C_3=C_4=0$。两端皆位置控制，等效“虚拟弹簧”耦合。$h_{11},h_{22}\neq0$、$h_{12}h_{21}\neq-1$——\emph{不透明}，但天然稳定、不需力传感器、对延迟鲁棒。代价：刚墙接触时残余弹性感（“软墙”），自由空间有残余阻尼。模仿学习常用的 ALOHA 平台（leader--follower 关节级“傀儡”同步、无任何力传感器与力反馈）是更退化的特例：仅 $C_1$ 单向位置映射、连 $C_2$ 力/位反馈都没有，故其实是\emph{单边}遥操作而非完整双边 PP。
  \item \textbf{PF（位置$\to$slave，力$\leftarrow$master）}：用 $C_1,C_4$。master 导纳控制（收 $f_e$ 算期望速度）、slave 位置控制，仅 slave 端需力传感器。$h_{11}\approx Z_{cm}$（master 自身阻抗，主要透明度损失源）、$h_{22}\approx0$。工业最常用的有力反馈架构，透明度--稳定性平衡合理。值得澄清一个常见误传：经典的 da Vinci 手术机器人（Si/X/Xi 各代）其实\emph{不向术者反馈力}（$C_4=0$），术者靠\emph{视觉}估计组织受力——即在力通道上是单边的、更接近 PP/单向位置映射，而非 PF；直到 2024 年获批的 da Vinci 5 才首次引入力反馈（Force Feedback，可测并让术者感到推拉力，官方称为该模态“首创”）。这一历史恰好例证了“安全/稳定优先于透明度”的工程取舍：在没有可靠力反馈通道的二十余年里，宁可牺牲力透明也要保稳。
  \item \textbf{FP（力$\to$slave，位置$\leftarrow$master）}：用 $C_3,C_2$。实践中几乎不用——slave 端力控在不确定接触下易失稳，且需 master 端力传感器（成本高、收益低）。
  \item \textbf{FF（力--力）}：四通道里两端皆力控，理论上最透明（$h_{11}\approx0,h_{22}\approx0$）。但工程上极难：力传感器噪声直入两端、力控双积分动力学低频不稳（位置漂移）、需两端力传感器、四通道增益高维耦合难调。多见于研究实验室。
\end{itemize}

\begin{table}[H]\centering\small
\caption{四种子架构的结构对比（透明度/稳定性为定性趋势）}
\label{tab:tt-arch}
\begin{tabular}{lcccll}
\toprule
架构 & $h_{11}$ & $h_{22}$ & 透明度 & 稳定性 & 力传感器 / 典型应用 \\
\midrule
PP & $\neq0$ & $\neq0$ & 低 & 高 & 0 个 / 学习数据采集 \\
PF & $\neq0$ & $\approx0$ & 中高 & 中 & 1 个(slave) / 工业力反馈、da Vinci~5 \\
FP & $\approx0$ & $\neq0$ & 中 & 低 & 1 个(master) / 几乎不用 \\
FF & $\approx0$ & $\approx0$ & 最高 & 最低 & 2 个 / 研究实验室 \\
\bottomrule
\end{tabular}
\end{table}

\begin{insight}[通道不是越多越好]\label{ins:tt-channels}
理想四通道全开可逼近 $\mathbf H_{\rm ideal}$，但需精确匹配多个控制器参数（Lawrence 给出 $C_1=Z_{cs},C_3=Z_{cm},C_2=1+C_5,C_4=1+C_6$ 一类条件\cite{lawrence1993stability}），任一失配都可能失稳。两通道（PP/PF）透明度有限却鲁棒得多。工业系统偏好 PF，因为可靠性优先于透明度——\emph{通道数越多越难调}。硬件接口也偏置架构选择：力矩接口（直驱关节）天然阻抗因果、位置接口（工业控制器）天然导纳因果，Hashtrudi-Zaad 与 Salcudean（2001）据此系统分析了阻抗/导纳 master $\times$ slave 的四种组合及选型原则\cite{hashtrudizaad2001analysis}。
\end{insight}

% ════════════════════════════════════════════════════════════════
\section{桥接：从透明度到绝对稳定性}
\label{sec:tt-bridge}

\paragraph{与 O 章 / P 章二端口骨架的对齐。}
本章不是凭空起灶。\cref{ch:operational-space} 的 \cref{sec:os-twoport} 已把“机器人 + 控制器”抽象为二端口、写出 $\mathbf Z/\mathbf Y/\mathbf H$ 三矩阵（\cref{eq:os-twoport}），并把所有力控范式统一为“端口 2 行为的选择”（位置控制 = 无穷输出阻抗的速度源、纯力控 = 零阻抗的力源、阻抗控制 = 介于其间的可编程连续谱）。\cref{ch:force-wbc} 的 \cref{sec:fw-twoport} 进一步把四种力控\emph{任务}归为二端口的不同工作点。本章把这套语言推进到\emph{人在环}：端口 1 不再接抽象控制器，而是接\emph{真实操作者}（带自身阻抗 $Z_h$），于是 $\mathbf H_{\rm ideal}$、$Z_{to}$ 这些“以操作者为参照”的量才有意义。三者的关系可一句话概括：O/P 章是\emph{机器人单边}的二端口（机器人 vs 环境），本章是\emph{双边 + 人}的二端口（人 vs 机器人 vs 环境）。

\paragraph{透明与稳定，在 $\mathbf H_{\rm ideal}$ 处正面相撞。}
本章建立了透明度这一侧的全部工具。下一章（\cref{ch:teleop-stability}）将建立稳定性这一侧的精确边界——Llewellyn 绝对稳定性判据，它把“对任意被动操作者 $Z_h$ 与任意被动环境 $Z_e$ 都闭环稳定”翻译成 $h_{11},h_{22},h_{12}h_{21}$ 上的四个充要条件，并定义稳定裕度 $\eta(\omega)\ge0$。此处只预告那个戏剧性的结论，作为承上启下的钩子：

\begin{insight}[零裕度相撞：透明度的代价（预告）]\label{ins:tt-tradeoff}
把理想透明的 $\mathbf H_{\rm ideal}$（$h_{11}=h_{22}=0,\ h_{12}h_{21}=-1$）代入绝对稳定性的裕度公式，可得 $\eta=0$——\emph{恰好}落在稳定与不稳定的分界线上，鲁棒裕度为零。这意味着理想透明系统名义上“刚好稳定”，但任何微小扰动（通信延迟引入的 $e^{-sT}$、电机惯量、传感器噪声、量化）都会把 $\eta$ 推向负、令系统失稳。要换得正裕度（$\eta>0$），唯一的代价是让 $h_{11},h_{22}$ 离零——物理上即在 master/slave 端\emph{加本地阻尼}，而这恰好让 $Z_{to}\neq Z_e$、操作者尝到额外的粘滞感。\textbf{透明度与稳定性是一对此消彼长的对手——这正是 \cref{ins:tt-core} 那句“根本矛盾”的精确数学形态。}完整推导（含 $\eta$ 的定义、四条件证明、延迟如何压低临界频率）见 \cref{ch:teleop-stability}。
\end{insight}

这把矛盾也指明了整个遥操作子部的逻辑链：既然加阻尼以保稳会牺牲透明、而延迟又会把临界频率压进人手操作频段（$1$--$10\,\mathrm{Hz}$），就需要\emph{不靠加阻尼}的办法在延迟下保被动——这正是 \cref{ch:teleop-wave}（波变量/散射变换，Anderson--Spong\cite{anderson1989bilateral}、Niemeyer--Slotine\cite{niemeyer1991stable}）与 \cref{ch:teleop-tdpa}（时域无源性控制 TDPA，Hannaford--Ryu\cite{hannaford2002time}）的使命；而运动/力的缩放映射与共享自主则留给 \cref{ch:teleop-mapping}。至于操作过程中机械臂与人/环境的物理碰撞安全，见 \cref{ch:arm-collision-safety}。

% ════════════════════════════════════════════════════════════════
\section{本章小结}
\label{sec:tt-summary}

本章把“人--主臂--从臂--环境”这条物理链抽象为一个二端口黑盒，用混合参数 $\mathbf H$ 量化它的透明度，并揭示透明与稳定的结构性矛盾。主线四步：被动符号约定下的端口变量（\cref{eq:tt-port-power}）$\to$ 因果匹配的混合参数 $\mathbf H$（\cref{eq:tt-hdef}）$\to$ 操作者感受阻抗 $Z_{to}$（\cref{eq:tt-zto}）$\to$ 理想透明 $\mathbf H_{\rm ideal}$（\cref{eq:tt-hideal}）。结论是：理想透明要求 master/slave“自身存在”消失，物理上不可达；逼近它又与稳定性零裕度相撞——后者由 \cref{ch:teleop-stability} 精确化。

\begin{table}[H]\centering\small
\caption{符号表}
\label{tab:tt-symbols}
\begin{tabular}{lll}
\toprule
符号 & 含义 & 首次出现 \\
\midrule
$f_h,v_m$ & 操作者施于 master 的力 / master 速度（端口 1） & \cref{fig:tt-twoport} \\
$f_e,v_s$ & 环境施于 slave 的力 / slave 速度（端口 2） & \cref{fig:tt-twoport} \\
$\nu_1,\nu_2$ & 流入网络的流变量 $\nu_1=v_m,\ \nu_2=-v_s$ & \cref{eq:tt-port-power} \\
$Z(s),Y(s)$ & 机械阻抗 $f/v$ / 导纳 $v/f$ & \cref{tab:tt-analogy} \\
$Z_e,Z_h$ & 环境阻抗 / 操作者阻抗 & \cref{thm:tt-zto} \\
$\mathbf H,h_{ij}$ & 混合参数矩阵及其元素 & \cref{eq:tt-hdef} \\
$\mathbf Z,\mathbf Y,\mathbf T,\mathbf S$ & 阻抗 / 导纳 / 传递 / 散射参数矩阵 & \cref{eq:tt-z2h,eq:tt-tdef} \\
$\Delta_h$ & $\det\mathbf H=h_{11}h_{22}-h_{12}h_{21}$ & \cref{eq:tt-zto} \\
$Z_{to}$ & 操作者感受阻抗 $f_h/v_m$ & \cref{eq:tt-zto} \\
$\mathbf H_{\rm ideal}$ & 理想透明的混合参数矩阵 & \cref{eq:tt-hideal} \\
$C_1$--$C_4$ & 四通道控制器（位置/力，前馈/反馈） & \cref{fig:tt-fourchan} \\
$\eta(\omega)$ & 绝对稳定裕度（详见 \cref{ch:teleop-stability}） & \cref{ins:tt-tradeoff} \\
\bottomrule
\end{tabular}
\end{table}

\begin{table}[H]\centering\small
\caption{定理 / 公式速查表}
\label{tab:tt-cheatsheet}
\begin{tabular}{lll}
\toprule
定理 / 公式 & 一句话说明 & 对应 \\
\midrule
端口功率 \cref{eq:tt-port-power} & $P=f_hv_m-f_ev_s$，承 \cref{eq:os-port-power} & \cref{sec:tt-port} \\
混合参数 \cref{eq:tt-hdef} & 因果匹配：master 力源、slave 运动源 & \cref{def:tt-hmatrix} \\
$\mathbf Z\to\mathbf H$ \cref{eq:tt-z2h} & 同一系统的无损换参 & \cref{sec:tt-convert} \\
级联 $\mathbf T_{\rm tot}=\mathbf T_m\mathbf T_{\rm ch}\mathbf T_s$ \cref{eq:tt-Ttot} & 传递矩阵的“相乘即串联” & \cref{prop:tt-cascade} \\
感受阻抗 \cref{eq:tt-zto} & $Z_{to}=h_{11}-h_{12}h_{21}Z_e/(1+h_{22}Z_e)$ & \cref{thm:tt-zto} \\
理想透明 \cref{eq:tt-hideal} & $\mathbf H_{\rm ideal}=[0,1;-1,0]$，幂系数比对 & \cref{thm:tt-hideal} \\
零裕度相撞 & $\mathbf H_{\rm ideal}\Rightarrow\eta=0$，透明 vs 稳定 & \cref{ins:tt-tradeoff} \\
\bottomrule
\end{tabular}
\end{table}

\begin{table}[H]\centering\small
\caption{知识点总表}
\label{tab:tt-knowledge}
\begin{tabular}{clll}
\toprule
\# & 知识点 & 核心要点 & 难度 \\
\midrule
1 & 核心矛盾 & 透明 vs 稳定是物理性权衡，非算法缺陷 & \(\bigstar\bigstar\) \\
2 & 端口变量/被动约定 & $\nu_2=-v_s$；$P=f_hv_m-f_ev_s$ & \(\bigstar\bigstar\) \\
3 & 混合参数 $\mathbf H$ & 因果匹配 master 力源/slave 运动源 & \(\bigstar\bigstar\bigstar\) \\
4 & 参数化转换 & $\mathbf Z\to\mathbf H$、$\mathbf T\to\mathbf H$、级联 & \(\bigstar\bigstar\bigstar\) \\
5 & 感受阻抗 $Z_{to}$ & 透明度的可计算度量；三特例 & \(\bigstar\bigstar\bigstar\) \\
6 & 理想透明 $\mathbf H_{\rm ideal}$ & 幂系数比对；$h_{11}=h_{22}=0$ & \(\bigstar\bigstar\bigstar\) \\
7 & 四子架构 & PP/PF/FP/FF 的透明度极限 & \(\bigstar\bigstar\bigstar\) \\
8 & 零裕度相撞（预告） & $\mathbf H_{\rm ideal}\Rightarrow\eta=0$ & \(\bigstar\bigstar\bigstar\bigstar\) \\
\bottomrule
\end{tabular}
\end{table}

\section*{练习}
\begin{practice}[练习]
\begin{enumerate}
  \item \textbf{[A 型 — $\mathbf H$ 参数计算]} 1-DOF 系统：master（质量 $M_m=0.5\,\mathrm{kg}$、阻尼 $B_m=2\,\mathrm{Ns/m}$）经 PD 控制耦合 slave（$M_s=1\,\mathrm{kg}$、$B_s=1\,\mathrm{Ns/m}$），位置跟踪 + 力反射。写出 $h_{11},h_{12},h_{21},h_{22}$ 的传递函数；取弹簧环境 $Z_e=K/s$，用 \cref{eq:tt-zto} 写出 $Z_{to}(j\omega)$ 并定性画出它与 $Z_e$ 的偏差。
  \item \textbf{[思考题]} 为什么 Lawrence（1993）用阻抗 $Z_{to}$ 而非力跟踪误差 $|f_h-f_e|$ 定义透明度？提示：考虑缩放系统 $f_h=10f_e,\ v_s=v_m/10$——力误差很大，但操作者的阻抗感受如何？（回看 \cref{sec:tt-zto} 的陷阱专栏。）
  \item \textbf{[B 型 — 换参验证]} 用 \cref{eq:tt-z2h} 把一个对称无源二端口（$z_{12}=z_{21}$）转换到 $\mathbf H$，验证所得 $\mathbf H$ 是否仍满足某种对称性。再用 \cref{eq:tt-t2h} 从一个理想变压器（$\mathbf T=\mathrm{diag}(n,1/n)$）求 $\mathbf H$，解释结果的物理含义。
  \item \textbf{[综合题]} 一台仅有位置编码器、无力传感器、$50\,\mathrm{Hz}$ 控制的双臂采集平台，能实现 PP/PF/FP/FF 中的哪些架构？不能实现的，分别缺什么硬件/能力？（提示：从 \cref{tab:tt-arch} 的“力传感器”列与因果性入手。）
\end{enumerate}
\end{practice}

\section*{延伸阅读}
\begin{itemize}
  \item \textbf{二端口与四通道奠基}：Hannaford\cite{hannaford1989design}（混合二端口模型与四通道架构的原始论文）；Lawrence\cite{lawrence1993stability}（透明度的严格定义、把 Llewellyn 判据引入遥操作、理想四通道条件）。
  \item \textbf{延迟下的无源通信}：Anderson--Spong\cite{anderson1989bilateral}（散射变换 $\to$ 任意常数时延无源）、Niemeyer--Slotine\cite{niemeyer1991stable}（波变量与传输线物理直觉）——是 \cref{ch:teleop-wave} 的源头。
  \item \textbf{时域无源性}：Hannaford--Ryu\cite{hannaford2002time}（TDPA/能量观测器，不依赖 LTI 假设）——是 \cref{ch:teleop-tdpa} 的源头；无源性遥操作的系统综述见 Nuño 等\cite{nuno2011passivity}。
  \item \textbf{选型框架}：Hashtrudi-Zaad--Salcudean\cite{hashtrudizaad2001analysis}（阻抗/导纳 master $\times$ slave 的四组合选型原则）。
  \item \textbf{历史与全景}：Ferrell\cite{ferrell1965remote}（时延实验）、Goertz\cite{goertz1954mechanical}（机械主从臂）；本章 \cref{sec:tt-bridge} 与 \cref{ch:operational-space}（\cref{sec:os-twoport}）、\cref{ch:force-wbc}（\cref{sec:fw-twoport}）的二端口骨架互链。
\end{itemize}
